% This is LLNCS.DEM the demonstration file of
% the LaTeX macro package from Springer-Verlag
% for Lecture Notes in Computer Science, version 1.1

\documentclass{llncs} 



%\documentclass[smallcondensed,natbib,10pt]{svjour3}

%\documentstyle{llncs}
%
\usepackage{graphicx}

\usepackage{multirow}
\usepackage{rotating}
%\usepackage[none]{hyphenat}

%\usepackage[utf8]{inputenc}
%\usepackage[turkish]{babel}

  \usepackage[turkish]{babel}
  \usepackage[utf8]{inputenc}
%  \usepackage[latin5]{inputenc}
  \usepackage[T1]{fontenc}
\setcounter{secnumdepth}{3}
 
\usepackage{enumerate}

% Citation styles
%\usepackage[natbibapa]{apacite}

%\usepackage{apacite}
%\usepackage[numbers]{apacite}

\begin{document}
%\usepackage[pdftex]{graphicx}
%%and Paul Clarke\inst{1}

\title{Yazılım Süreç Değişimlerinin ulusal yazılım endüstrisine olan etkileri: Bir vaka çalışması}

\subtitle{Exploring the implication of software processes change in national software industry: A case study}



%\subtitle{Endüstriyel vaka çalışması}
%\subtitle{A proposal for the research into Game Theory}
\author{Murat Yilmaz\inst{1}}
\institute{Çankaya Üniversitesi, Bilgisayar Mühendisliği, Ankara, Türkiye
}



\maketitle

\begin{abstract}


\end{abstract}

\begin{keywords}
Yazılım yönetimi, Sosyal etmenler, İşe alıştırma, oryantasyon
\end{keywords}



\section{Giriş}



aaaaa~\cite{yilmaz2013software}


\section{Araştırmanın Temelleri}


\subsection{İşe Alıştırma (Onboarding) Kavramı}




\section{Endüstiyel Vaka Çalışması}




Odak grup çalışması esnasında, araştırmacılara aşağıda belirtilen örnek sorular yönlendirilerek katılımcıların konu üzerinde tartışmaları sağlanmıştır.

%\begin{itemize}
%\item Yazılım geliştirme süreçlerinde işe alıştırma ile ilgili problemler nelerdir?



%\textbf{K2} \textquotedblleft{}\textquotedblright{}

%\textit{\begin{quote}
%\textbf{Söyleşi ayrıntısı:}\textquotedblleft{} \textquotedblright{}
%\end{quote}} 

%\textit{\begin{quote}
%\textbf{Söyleşi ayrıntısı:} \textquotedblleft{}Onboarding (işe alıştırma), mevcut ekip üyelerinin yeni gelen bir kişiye takım ile çalışmaya alışmalarına yardımcı olmak için gereken süreci denetler. Yeni gelenlerin bizim projelerde büyük miktarda sosyal ve teknik bilgiyi almaları beklenir; bizim yazılım geliştirme takımı için, bu genellikle üzerinde çalışacakları kaynak kodun anlaşılmasını içerir; bu nedenle, geliştirilecek ciddi oyun, \textbf{program anlama örnekleri} sunmalıdır. Ayrıca, modüler olmalı ve yönetici arayüzü ile bilgi girişi sağlanmalıdır.\textquotedblright{}
%\end{quote}} 


%Diğer bir katılımcı işe alıştırmanın takım genişletmek için zorunlu olduğunu söyleyerek şu bilgileri iletmiştir:



%Birçok katılımcı aşağıdaki önermelerin doğru olduğunu düşünmüştür.


%\newpage




\section{Sonuçlar ve Gelecek Çalışmalar}



%(bkz. Şekil~\ref{fig:dem2}). 

%\begin{figure}[h!tb]
% \begin{center}
% \shorthandoff{=}
 %\raggedleft    
%    \includegraphics[scale=0.15]{yontem.png}
%      \shorthandon{=}
%    \caption{Araştırmada Kullanılan Yöntem Adımları.}\label{fig:dem2}
%  \end{center}
%\end{figure}




%\section*{Teşekkür}

%Bu çalışma İrlanda Bilimsel ve Teknolojik Araştırma Kurumunun, Lero – İrlanda Yazılım Mühendisliği Araştırma Merkezi’ne verdiği $10{\rm /CE}{\rm /I1} 855$  kodlu destek programı dahilinde desteklenmektedir.
%splncs.bst
%\bibliographystyle{splncs03.bst}
%\bibliographystyle{splncs.bst}

%\bibliographystyle{apacite}
\bibliographystyle{splncs}
%\bibliography{bibfilename} 

% BibTeX users please use one of
%\bibliographystyle{spbasic}      % basic style, author-year citations
%\bibliographystyle{spmpsci}      % mathematics and physical sciences
%\bibliographystyle{spphys}       % APS-like style for physics
%\bibliography{}   % name your BibTeX data base

\bibliography{uludag}


%\bibliographystyle{plain}
%\bibliography{uyms2017}
\end{document}

